%% GENERATED by tools/gen_latex.py -- edit content/agents03.py instead.
%% Build:  latexmk -pdf agents03.tex     (or: pdflatex agents03.tex, twice)
\documentclass[aspectratio=169,11pt,t]{beamer}
\usepackage{itbpro}

\coursetitle{Model penalar, dan kapan penalaran layak dibayar}
\coursesubtitle{Memindahkan komputasi dari waktu latih ke waktu pakai --- apa yang dibelinya, berapa harganya, dan kenapa agen sudah melakukan sebagian dari itu.}
\coursekicker{Chapter 3}
\coursesource{Grootendorst \& Alammar, An Illustrated Guide to AI Agents --- bab 3}
\courseduration{3 jam (2 sesi)}
\coursepresenter{Hendri Karisma}{Instructor}
\coursebrand{AI for Professional \textperiodcentered\ ITB \textperiodcentered\ Ch 3}

\title{Model penalar, dan kapan penalaran layak dibayar}
\author{Hendri Karisma}
\date{}

\begin{document}
\covertitle

\framekicker{Learning outcomes}
\begin{frame}[shrink=25]{Session Objectives}
\leadpar{By the end of this session, participants can:}
\begin{isteps}
\item \textbf{Membedakan komputasi saat latih dari komputasi saat pakai}, dan menyebutkan mana yang bisa diatur setelah model dikirim.
\item \textbf{Menghitung apa yang dibeli oleh pengambilan N contoh}, dan menyebutkan asumsi yang membuat hitungan itu jadi batas atas.
\item \textbf{Menjelaskan kenapa pemeriksa mengalahkan penambahan contoh}, dan menyebutkan syarat sebuah tugas punya pemeriksa.
\item \textbf{Menyebutkan berapa token yang dibayar tapi tidak terlihat} pada model penalar, dan akibatnya pada agen berlangkah banyak.
\item \textbf{Memutuskan kapan TIDAK memakai model penalar} di dalam gelung agen, dengan alasan yang bisa diukur.
\item \textbf{Menyebutkan tiga cara penalaran panjang justru memperburuk} hasil.
\end{isteps}

\vskip 6pt
\vskip 4pt
\reslink{Site}{Course home}{https://hendrikarisma.my.id/rs/ai-products-course/index.html}
\reslink{Book}{Maarten Grootendorst \& Jay Alammar, An Illustrated Guide to AI Agents (O'Reilly Media, Early Release)}{https://www.oreilly.com/library/view/an-illustrated-guide/9798341662681/}
\vskip 2pt
\end{frame}

\coursesection{01}{Dua tempat menaruh komputasi}{Satu dibayar sekali di depan, satu dibayar tiap permintaan.}

\framekicker{Pergeseran}
\begin{frame}[shrink=25]{Yang berubah bukan modelnya, melainkan kapan usahanya dikeluarkan}
\begin{center}
\includegraphics[width=0.92\linewidth,height=0.60\textheight,keepaspectratio]{figs/agents03-compute.tex.pdf}
\end{center}
\vskip 3pt{\color{ink3}\fontsize{7.4}{9.4}\selectfont Bobot dibayar sekali; berpikir dibayar tiap kali.\par}
{\fontsize{9.4}{12.6}\selectfont\color{fg2}Selama bertahun-tahun, cara menaikkan kemampuan adalah \textbf{membesarkan pelatihannya}: lebih banyak data, lebih banyak parameter, lebih banyak GPU. Hasilnya bobot yang tetap --- pintar atau tidak, ia sama saja untuk tiap permintaan.\par}\vskip 4pt
\begin{iband}[signal]
Yang berubah: sebagian usaha dipindah ke \textbf{saat pemakaian}. Model yang sama diberi kesempatan berpikir lebih lama pada soal yang lebih sulit --- dan itu berarti biayanya jadi \hilite{sesuatu yang Anda atur per permintaan}, bukan sesuatu yang sudah tetap.
\end{iband}
\end{frame}

\framekicker{Pergeseran}
\begin{frame}[shrink=25]{Kenapa ini kabar baik bagi orang yang membangun sistem}
\begin{minipage}[t]{0.485\textwidth}
{\fontsize{9.4}{12.6}\selectfont\color{fg2}\textbf{Komputasi saat latih} adalah keputusan orang lain. Anda menerimanya sebagaimana adanya, dan satu-satunya kendali Anda adalah memilih model.\par}\vskip 4pt

\end{minipage}%
\hfill
\begin{minipage}[t]{0.485\textwidth}
{\fontsize{9.4}{12.6}\selectfont\color{fg2}\textbf{Komputasi saat pakai} adalah keputusan Anda, per permintaan. Soal mudah dijawab murah; soal sulit diberi anggaran lebih besar.\par}\vskip 4pt

\end{minipage}%
{\fontsize{9.4}{12.6}\selectfont\color{fg2}Ini pola yang sudah dikenal di tempat lain: cache yang murah untuk yang sering, jalur mahal untuk yang jarang. Yang baru adalah \textbf{kualitas jawaban} kini ikut jadi sesuatu yang bisa dibeli dengan uang di saat pemakaian.\par}\vskip 4pt
\begin{iband}[signal]
Dan kabar buruknya, yang muncul di sisa bab ini: yang dibeli \textbf{melandai}, sedangkan yang dibayar tidak.
\end{iband}
\end{frame}

\coursesection{02}{Cara termurah dulu}{Sebelum membayar model penalar, ada yang bisa dicoba dengan harga nol.}

\framekicker{Prompt}
\begin{frame}[shrink=25]{Meminta langkahnya adalah pengungkit yang paling murah}
{\fontsize{9.4}{12.6}\selectfont\color{fg2}Meminta model menuliskan langkah-langkahnya sebelum menjawab menaikkan ketepatan pada soal berlangkah banyak --- dan harganya hanya token keluaran tambahan.\par}\vskip 4pt
{\fontsize{9.4}{12.6}\selectfont\color{fg2}Alasannya masuk akal begitu diingat bahwa model menghasilkan satu token dari semua token sebelumnya: \textbf{langkah yang ditulis menjadi bagian dari konteks untuk langkah berikutnya.} Ia bukan penjelasan setelah keputusan --- ia bagian dari perhitungannya.\par}\vskip 4pt
\begin{iband}[signal]
Karena itu penjelasan yang diminta \textbf{sesudah} jawaban keluar adalah hal yang berbeda sama sekali: itu karangan tentang keputusan yang sudah diambil, dan \hilite{tidak boleh dibaca sebagai alasan sebenarnya}
\end{iband}
\end{frame}

\framekicker{Prompt}
\begin{frame}[shrink=25]{Di dalam agen, sebagian besar ini sudah terjadi}
\begin{center}

\includegraphics[width=0.92\linewidth,height=0.60\textheight,keepaspectratio]{figs/agents03-where.tex.pdf}
\end{center}
\vskip 3pt{\color{ink3}\fontsize{7.4}{9.4}\selectfont Dua tempat penalaran bisa terjadi, dan hanya satu yang bisa dilihat.\par}
{\fontsize{9.4}{12.6}\selectfont\color{fg2}Gelung agen \textbf{sudah} memecah masalah jadi langkah. Bedanya dengan penalaran di dalam model bukan ada atau tidaknya langkah, melainkan \textbf{apakah langkah itu bisa dilihat, diukur, dan dihentikan} --- dan akibatnya, menaruh model penalar di dalam gelung kadang berarti membayar dua kali untuk hal yang sama.\par}\vskip 4pt
\end{frame}

\coursesection{03}{Bertanya berkali-kali}{Dan berapa persisnya itu membeli sesuatu.}

\framekicker{Banyak contoh}
\begin{frame}[shrink=25]{Ambil N jawaban, pakai yang terbanyak}
{\fontsize{9.4}{12.6}\selectfont\color{fg2}Kalau satu jawaban bisa salah, ambil beberapa dan pilih yang paling sering muncul. Gagasannya tua dan masuk akal: kesalahan acak cenderung berbeda-beda, jawaban benar cenderung sama.\par}\vskip 4pt
{\fontsize{9.4}{12.6}\selectfont\color{fg2}Yang jarang dihitung orang adalah \textbf{berapa banyak} itu membeli. Untungnya itu bukan soal selera --- kalau tiap contoh benar dengan peluang \textit{p} dan contohnya saling bebas, peluang suara terbanyak benar adalah jumlah suku binomial di atas N/2. Bisa dihitung persis.\par}\vskip 4pt
\end{frame}

\framekicker{Banyak contoh}
\begin{frame}[shrink=25]{Dihitung, bukan ditaksir --- dan asumsinya dicetak}
\begin{center}
\includegraphics[width=0.97\linewidth,height=0.56\textheight,keepaspectratio]{figs/agents03-vote.tex.pdf}
\end{center}
\vskip 3pt{\color{ink3}\fontsize{7.4}{9.4}\selectfont Peluang benar setelah suara terbanyak dari N contoh, dihitung dari sebaran binomial. Langkahi: satu contoh, sampai lima, lalu semuanya.\par}
{\fontsize{9.4}{12.6}\selectfont\color{fg2}Dari 1 ke 5 contoh: \textbf{+8,3 poin} dengan biaya 5$\times$. Dari 5 ke 21: \textbf{+14,3 poin} dengan biaya 21$\times$. Kenaikannya melandai, biayanya tidak --- dan karena contohnya berkorelasi, yang benar-benar didapat selalu lebih kecil dari angka itu.\par}\vskip 4pt
\end{frame}

\framekicker{Banyak contoh}
\begin{frame}[shrink=25]{Satu syarat yang membuatnya bisa dipakai sama sekali}
{\fontsize{9.4}{12.6}\selectfont\color{fg2}Suara terbanyak hanya berarti kalau jawabannya \textbf{bisa dibandingkan}. Untuk angka atau label, mudah. Untuk paragraf, dua jawaban benar bisa berbeda kata sepenuhnya dan tidak ada suara yang bisa dihitung.\par}\vskip 4pt
\begin{minipage}[t]{0.3253\textwidth}
\begin{icardgood}
\cardh{Cocok}
Angka, klasifikasi, pilihan alat, keputusan ya/tidak.
\end{icardgood}
\end{minipage}%
\hfill
\begin{minipage}[t]{0.3253\textwidth}
\begin{icard}
\cardh{Perlu usaha}
Keluaran terstruktur --- bandingkan medannya, bukan teksnya.
\end{icard}
\end{minipage}%
\hfill
\begin{minipage}[t]{0.3253\textwidth}
\begin{icardbad}
\cardh{Tidak cocok}
Prosa bebas. Suara terbanyak di sini hampir selalu berarti membandingkan gaya, bukan kebenaran.
\end{icardbad}
\end{minipage}%
\begin{iband}[signal]
Di dalam agen, yang paling sering pantas divoting justru bukan jawabannya, melainkan \textbf{pilihan alat dan argumennya} --- dan itu kebetulan hal yang paling mudah dibandingkan.
\end{iband}
\end{frame}

\framekicker{Pemeriksa}
\begin{frame}[shrink=25]{Pemeriksa mengalahkan penambahan contoh, kalau ada}
\begin{center}
\includegraphics[width=0.92\linewidth,height=0.60\textheight,keepaspectratio]{figs/agents03-verifier.tex.pdf}
\end{center}
\vskip 3pt{\color{ink3}\fontsize{7.4}{9.4}\selectfont Menghasilkan banyak calon jadi berguna begitu ada yang bisa menolak.\par}
{\fontsize{9.4}{12.6}\selectfont\color{fg2}Suara terbanyak menebak mana yang benar dari \textbf{kesepakatan}. Pemeriksa mengetahuinya dari \textbf{pengujian}. Kalau tugasnya punya pemeriksa --- kode yang dijalankan, hitungan yang dicocokkan, skema yang divalidasi, kebijakan yang diperiksa --- maka menghasilkan lima calon dan menolak yang gagal jauh lebih kuat daripada memilih yang paling populer.\par}\vskip 4pt
\begin{iband}[signal]
Syaratnya satu, dan itu syarat yang tegas: \textbf{memeriksa harus lebih mudah daripada menjawab.} Kalau memeriksanya sama sulitnya, Anda baru saja menggandakan masalahnya.
\end{iband}
\end{frame}

\framekicker{Pemeriksa}
\begin{frame}[shrink=25]{Apa yang bisa jadi pemeriksa dalam sistem sungguhan}
\begin{itable}{>{\raggedright\arraybackslash}p{0.2444\textwidth}>{\raggedright\arraybackslash}p{0.4136\textwidth}>{\raggedright\arraybackslash}p{0.282\textwidth}}
\thead{Pemeriksa} & \thead{Menolak apa} & \thead{Ongkosnya} \\ \midrule
Menjalankan kodenya & Kode yang tidak jalan atau gagal uji & Sandbox --- dan itu bukan hal sepele \\
Validasi skema & Argumen alat yang salah bentuk atau tipe & Hampir nol, dan wajib ada \\
Hitung ulang dengan kode & Aritmetika yang dikarang model & Nol, dan mengejutkan jarang dipakai \\
Cocokkan dengan hasil alat & Angka yang tidak ada asal-usulnya & Rendah --- butuh jejak yang rapi \\
Aturan kebijakan sebagai kode & Keputusan yang melanggar syarat & Sedang --- tapi harus ada juga untuk audit \\
\end{itable}
{\fontsize{9.4}{12.6}\selectfont\color{fg2}Perhatikan bahwa empat dari lima \textbf{bukan model}. Pemeriksa terbaik biasanya kode biasa, dan itu sebabnya bab ini berakhir pada nasihat yang membosankan: tulis pemeriksanya dulu.\par}\vskip 4pt
\end{frame}

\framekicker{Hukum skala}
\begin{frame}[shrink=25]{Kenapa membesarkan pelatihan berhenti jadi jawaban tunggal}
{\fontsize{9.4}{12.6}\selectfont\color{fg2}Selama satu dekade, hubungannya rapi: lebih banyak data dan parameter menghasilkan model yang lebih baik, dengan cara yang bisa diramalkan. Itu yang membuat perlombaan ukuran masuk akal.\par}\vskip 4pt
\begin{minipage}[t]{0.3253\textwidth}
\begin{icard}
\cardh{Hasilnya melandai}
Kenaikan yang sama menuntut lipatan sumber daya yang makin besar.
\end{icard}
\end{minipage}%
\hfill
\begin{minipage}[t]{0.3253\textwidth}
\begin{icard}
\cardh{Datanya terbatas}
Teks berkualitas tinggi bukan sumber daya tak terbatas, dan bagian termudahnya sudah dipakai.
\end{icard}
\end{minipage}%
\hfill
\begin{minipage}[t]{0.3253\textwidth}
\begin{icard}
\cardh{Ongkosnya di depan}
Salah taruhan pada pelatihan berarti berbulan-bulan dan biaya besar yang tidak bisa ditarik.
\end{icard}
\end{minipage}%
{\fontsize{9.4}{12.6}\selectfont\color{fg2}Komputasi saat pakai tidak menggantikan itu --- ia menambah \textbf{sumbu kedua}. Dan sumbu kedua ini punya sifat yang menarik bagi orang yang membangun produk: ia bisa diatur sesudah semuanya dikirim.\par}\vskip 4pt
\end{frame}

\framekicker{Hukum skala}
\begin{frame}[shrink=25]{Dua sumbu, dan yang kedua milik Anda}
\begin{itable}{>{\raggedright\arraybackslash}p{0.2068\textwidth}>{\raggedright\arraybackslash}p{0.3666\textwidth}>{\raggedright\arraybackslash}p{0.3666\textwidth}}
\thead{} & \thead{Komputasi saat latih} & \thead{Komputasi saat pakai} \\ \midrule
Dibayar & Sekali, di depan & Tiap permintaan \\
Diputuskan oleh & Penyedia model & \textbf{Anda}, per permintaan \\
Bisa diubah & Tidak, tanpa melatih ulang & Ya, dengan satu parameter \\
Kalau salah & Berbulan-bulan hilang & Satu permintaan lebih mahal \\
Batasnya & Data dan uang & Kesabaran pengguna dan tagihan \\
\end{itable}
\begin{iband}[signal]
Untuk pembangun agen, baris ketiga yang paling penting: \textbf{ini satu-satunya tuas kualitas yang Anda pegang sendiri} tanpa mengganti model.
\end{iband}
\end{frame}

\framekicker{Melatihnya}
\begin{frame}[shrink=25]{Cara pertama: tunjukkan contoh penalaran yang bagus}
{\fontsize{9.4}{12.6}\selectfont\color{fg2}Kalau model dilatih pada kumpulan contoh yang jawabannya disertai langkah-langkah, ia belajar menghasilkan langkah itu sendiri --- tanpa diminta lewat prompt.\par}\vskip 4pt
{\fontsize{9.4}{12.6}\selectfont\color{fg2}Yang mengejutkan dari arah penelitian ini adalah \textbf{berapa sedikit contoh yang dibutuhkan}. Ribuan contoh penalaran yang dipilih dengan cermat bisa mengalahkan ratusan ribu contoh yang dikumpulkan asal --- kualitas dan keberagaman mengalahkan jumlah.\par}\vskip 4pt
\begin{iband}[signal]
Pelajaran yang berlaku juga untuk kumpulan uji Anda sendiri: \hilite{dua puluh kasus yang dipilih baik lebih berguna daripada dua ribu kasus yang dikumpulkan sembarangan}
\end{iband}
\end{frame}

\framekicker{Melatihnya}
\begin{frame}[shrink=25]{Cara kedua: beri imbalan pada jawaban yang benar, bukan pada caranya}
{\fontsize{9.4}{12.6}\selectfont\color{fg2}Pendekatan yang lebih mengejutkan: jangan tunjukkan cara menalar sama sekali. Beri saja soal yang jawabannya bisa \textbf{diperiksa otomatis}, biarkan model mencoba, dan beri imbalan ketika jawabannya benar.\par}\vskip 4pt
\begin{isteps}
\item \{'h': 'Yang membuatnya bekerja adalah pemeriksanya', 'p': 'Matematika dan pemrograman punya jawaban yang bisa dicek mesin. Di situlah metode ini paling berhasil, dan bukan kebetulan.'\}
\item \{'h': 'Perilaku memeriksa ulang muncul sendiri', 'p': 'Model mulai mengoreksi dirinya di tengah jalan tanpa pernah diberi contoh cara melakukannya --- sebab itu menaikkan peluang jawabannya benar.'\}
\item \{'h': 'Panjang penalarannya ikut tumbuh sendiri', 'p': 'Bukan karena ditargetkan, melainkan karena berpikir lebih lama kebetulan berkorelasi dengan benar pada soal-soal itu.'\}
\end{isteps}
\begin{iband}[signal]
Batasnya jadi jelas begitu polanya dilihat: metode ini \textbf{butuh domain yang bisa diperiksa mesin.} Untuk tugas yang benarnya subjektif, tidak ada imbalan yang bisa dihitung.
\end{iband}
\end{frame}

\framekicker{Melatihnya}
\begin{frame}[shrink=25]{Kenapa dua tahap, bukan satu}
\begin{minipage}[t]{0.485\textwidth}
{\fontsize{9.4}{12.6}\selectfont\color{fg2}\textbf{Imbalan saja} menghasilkan penalaran yang efektif tetapi sering kacau dibaca: berpindah bahasa di tengah, mengulang, atau menulis dalam bentuk yang tidak berguna bagi manusia.\par}\vskip 4pt
{\fontsize{9.4}{12.6}\selectfont\color{fg2}Wajar --- tidak ada dalam imbalannya yang meminta supaya bisa dibaca.\par}\vskip 4pt

\end{minipage}%
\hfill
\begin{minipage}[t]{0.485\textwidth}
{\fontsize{9.4}{12.6}\selectfont\color{fg2}\textbf{Contoh dulu, baru imbalan} menghasilkan penalaran yang rapi dan tetap efektif: contoh mengajarkan bentuknya, imbalan mengajarkan mana yang berhasil.\par}\vskip 4pt
{\fontsize{9.4}{12.6}\selectfont\color{fg2}Ini pola yang sama dengan bab 2: SFT mengajarkan bentuk, RL mengajarkan pilihan.\par}\vskip 4pt

\end{minipage}%
\begin{iband}[signal]
Bagi pemakai, akibatnya praktis: \textbf{jejak penalaran yang bisa dibaca bukan sifat alami model} --- ia hasil tahap pelatihan yang sengaja ditambahkan. Jangan menganggapnya jaminan.
\end{iband}
\end{frame}

\framekicker{Melatihnya}
\begin{frame}[shrink=25]{Memaksa berpikir lebih lama, dari luar}
{\fontsize{9.4}{12.6}\selectfont\color{fg2}Ada trik yang sangat sederhana dan bekerja mengejutkan baik: ketika model hendak berhenti berpikir, \textbf{jangan biarkan} --- sisipkan penanda yang membuatnya melanjutkan, beberapa kali.\par}\vskip 4pt
{\fontsize{9.4}{12.6}\selectfont\color{fg2}Kebalikannya juga bekerja: memotong penalaran pada batas token tertentu dan memaksa jawaban keluar. Dua-duanya menjadikan panjang berpikir sebagai \textbf{tombol yang bisa diputar saat pemakaian}, bukan sifat model.\par}\vskip 4pt
\begin{iband}[signal]
Inilah bentuk paling murni dari gagasan bab ini: satu bobot yang sama, kualitas yang berbeda-beda, ditentukan oleh berapa banyak komputasi yang Anda izinkan per permintaan.
\end{iband}
\end{frame}

\coursesection{04b}{Arah yang sedang dikerjakan}{Tiga hal yang belum selesai, dan pantas dilihat sebelum bertaruh padanya.}

\framekicker{Arah riset}
\begin{frame}[shrink=25]{Penalaran yang lebih hemat}
{\fontsize{9.4}{12.6}\selectfont\color{fg2}Masalah paling nyata dari model penalar adalah \textbf{borosnya}, dan sebagian besarnya terbuang: model sering terus berpikir pada soal yang sudah jelas jawabannya sejak awal.\par}\vskip 4pt
\begin{minipage}[t]{0.3253\textwidth}
\begin{icard}
\cardh{Berhenti kalau sudah yakin}
Menghentikan penalaran begitu jawabannya stabil, bukan setelah anggaran habis.
\end{icard}
\end{minipage}%
\hfill
\begin{minipage}[t]{0.3253\textwidth}
\begin{icard}
\cardh{Arahkan per permintaan}
Soal mudah ke model cepat, soal sulit ke model penalar --- keputusan dibuat sebelum menjawab.
\end{icard}
\end{minipage}%
\hfill
\begin{minipage}[t]{0.3253\textwidth}
\begin{icard}
\cardh{Penalaran yang lebih padat}
Melatih model menalar dengan token lebih sedikit untuk hasil yang sama.
\end{icard}
\end{minipage}%
{\fontsize{9.4}{12.6}\selectfont\color{fg2}Yang kedua bisa Anda lakukan \textbf{hari ini}, tanpa menunggu penelitian: sebuah pengklasifikasi kecil di depan gelung yang memilih model mana yang dipakai. Itu rekayasa biasa.\par}\vskip 4pt
\end{frame}

\framekicker{Arah riset}
\begin{frame}[shrink=25]{Menalar tanpa menuliskannya}
{\fontsize{9.4}{12.6}\selectfont\color{fg2}Menulis penalaran dalam bentuk kata memaksa model menyalurkan pikirannya lewat saluran yang sempit: satu token pada satu waktu, dan hanya yang bisa diucapkan.\par}\vskip 4pt
{\fontsize{9.4}{12.6}\selectfont\color{fg2}Arah penelitian yang sedang dikerjakan membiarkan model menalar di dalam \textbf{ruang vektornya sendiri} dan hanya mengeluarkan hasil akhirnya. Lebih hemat, dan mungkin lebih kuat.\par}\vskip 4pt
\begin{iband}[signal]
Tetapi perhatikan harganya bagi orang yang membangun sistem bertanggung jawab: \textbf{penalaran yang tidak berbentuk kata adalah penalaran yang tidak bisa dibaca, diaudit, atau dibantah.} Untuk keputusan yang diawasi regulator, itu bukan pertukaran yang netral.
\end{iband}
\end{frame}

\framekicker{Arah riset}
\begin{frame}[shrink=25]{Menalar atas gambar, bukan hanya atas teks}
{\fontsize{9.4}{12.6}\selectfont\color{fg2}Penalaran berlangkah pada masukan gambar masih jauh lebih lemah daripada pada teks. Model bisa menjelaskan isi sebuah gambar dengan baik, dan tetap gagal pada soal yang menuntut beberapa langkah \textbf{di atas} isi itu.\par}\vskip 4pt
{\fontsize{9.4}{12.6}\selectfont\color{fg2}Ini penting untuk agen yang bekerja di lapangan --- membaca faktur, layar, foto dokumen --- dan bab 9 kembali ke sana. Untuk sekarang cukup satu kesimpulan praktis.\par}\vskip 4pt
\begin{iband}[signal]
Kalau tugasnya menuntut penalaran di atas gambar, \textbf{ubah dulu gambarnya jadi data terstruktur dengan alat}, lalu nalarkan atas datanya. Jangan meminta satu panggilan mengerjakan keduanya.
\end{iband}
\end{frame}

\coursesection{04}{Model yang dilatih untuk menalar}{Dan tagihan yang tidak terlihat di layar.}

\framekicker{Model penalar}
\begin{frame}[shrink=25]{Penalaran yang dilatih, bukan yang diminta}
{\fontsize{9.4}{12.6}\selectfont\color{fg2}Alih-alih meminta model menuliskan langkahnya lewat prompt, model bisa \textbf{dilatih} untuk melakukannya sendiri --- diberi imbalan ketika jalur penalaran yang panjang berujung pada jawaban benar.\par}\vskip 4pt
{\fontsize{9.4}{12.6}\selectfont\color{fg2}Hasilnya model yang, diberi soal sulit, menghabiskan ratusan sampai ribuan token untuk dirinya sendiri sebelum mengeluarkan jawaban. Pada soal matematika, pemrograman, dan perencanaan berlangkah banyak, bedanya nyata.\par}\vskip 4pt
\begin{iband}[signal]
Dan satu temuan yang layak diingat karena berlawanan dengan dugaan: \textbf{perilaku memeriksa ulang dan mengoreksi diri bisa muncul dari pelatihan imbalan saja}, tanpa seorang pun menuliskan contoh cara memeriksa ulang.
\end{iband}
\end{frame}

\framekicker{Model penalar}
\begin{frame}[shrink=25]{Token yang Anda bayar dan tidak pernah Anda lihat}
{\fontsize{9.4}{12.6}\selectfont\color{fg2}Token penalaran adalah \textbf{token keluaran}. Ditagih seperti token keluaran, dihasilkan selambat token keluaran, dan pada banyak penyedia \textbf{tidak ditampilkan} kepada Anda.\par}\vskip 4pt
\begin{itable}{>{\raggedright\arraybackslash}p{0.3196\textwidth}>{\raggedright\arraybackslash}p{0.282\textwidth}>{\raggedright\arraybackslash}p{0.3384\textwidth}}
\thead{Satu tugas} & \thead{Token keluaran} & \thead{Yang terlihat} \\ \midrule
Model biasa & 300 & 300 --- semuanya \\
Model penalar & 2 800 & 300 --- sisanya \textbf{89\% tak terlihat} \\
Agen 6 langkah, model biasa & 1 800 & 1 800 \\
Agen 6 langkah, model penalar & \textbf{16 800} & 1 800 \\
\end{itable}
\begin{iband}[signal]
Asumsinya dicetak supaya bisa diganti: 300 token jawaban, 2 500 token berpikir per langkah. Yang tidak berubah adalah bentuknya --- \hilite{di dalam gelung agen, biaya penalaran dikalikan jumlah langkah}
\end{iband}
\end{frame}

\framekicker{Model penalar}
\begin{frame}[shrink=25]{Anggaran berpikir adalah tombol, dan pantas dipakai}
\begin{center}
\includegraphics[width=0.92\linewidth,height=0.60\textheight,keepaspectratio]{figs/agents03-budget.tex.pdf}
\end{center}
\vskip 3pt{\color{ink3}\fontsize{7.4}{9.4}\selectfont Tidak semua permintaan pantas mendapat anggaran berpikir yang sama.\par}
{\fontsize{9.4}{12.6}\selectfont\color{fg2}Sebagian besar penyedia model penalar menyediakan cara membatasi berapa banyak token boleh dipakai untuk berpikir. Memakainya bukan penghematan receh: pada agen, ia berlipat dengan jumlah langkah.\par}\vskip 4pt
\begin{iband}[signal]
Pola yang paling sering menang di produksi: \textbf{model cepat untuk memilih alat, model penalar hanya untuk langkah yang benar-benar sulit.} Memilih alat jarang butuh penalaran panjang; ia butuh kepatuhan pada skema.
\end{iband}
\end{frame}

\coursesection{05}{Ketika penalaran memperburuk}{Tiga cara, dan semuanya terlihat seperti keberhasilan sampai diperiksa.}

\framekicker{Batas}
\begin{frame}[shrink=25]{Berpikir lebih lama bukan berarti lebih benar}
\begin{center}
\includegraphics[width=0.92\linewidth,height=0.60\textheight,keepaspectratio]{figs/agents03-fail.tex.pdf}
\end{center}
\vskip 3pt{\color{ink3}\fontsize{7.4}{9.4}\selectfont Tiga jalan dari penalaran panjang menuju kesalahan yang meyakinkan.\par}
{\fontsize{9.4}{12.6}\selectfont\color{fg2}Penalaran panjang menambah \textbf{jumlah langkah yang bisa salah}. Satu kekeliruan di langkah ketiga akan dibangun di atasnya oleh sepuluh langkah berikutnya, dan hasil akhirnya terdengar jauh lebih meyakinkan daripada tebakan langsung --- sebab ia datang dengan alasan.\par}\vskip 4pt
\begin{iband}[signal]
Ini alasan lain kenapa penilaian harus melihat \textbf{jejaknya}: jawaban yang benar dari penalaran yang keliru dan jawaban yang benar dari penalaran yang benar terlihat sama di kolom hasil.
\end{iband}
\end{frame}

\framekicker{Batas}
\begin{frame}[shrink=25]{Tugas yang justru dirugikan oleh penalaran panjang}
\begin{minipage}[t]{0.494\textwidth}
\begin{icardbad}
\cardh{Pencarian fakta sederhana}
Jawabannya ada di hasil alat. Berpikir lama hanya menambah kesempatan mengarang penjelasan yang tidak diminta.
\end{icardbad}
\end{minipage}%
\hfill
\begin{minipage}[t]{0.494\textwidth}
\begin{icardbad}
\cardh{Memilih alat yang sudah jelas}
Kalau hanya satu alat yang masuk akal, penalaran panjang adalah biaya murni.
\end{icardbad}
\end{minipage}%


\vskip 4pt

\begin{minipage}[t]{0.494\textwidth}
\begin{icard}
\cardh{Apa pun yang berhadapan dengan pengguna}
Jeda beberapa detik per langkah $\times$ enam langkah adalah pengalaman yang buruk, betapapun bagus jawabannya.
\end{icard}
\end{minipage}%
\hfill
\begin{minipage}[t]{0.494\textwidth}
\begin{icard}
\cardh{Keluaran berformat ketat}
Penalaran panjang menaikkan peluang menyimpang dari format yang diminta, bukan menurunkannya.
\end{icard}
\end{minipage}%
\end{frame}

\framekicker{Batas}
\begin{frame}[shrink=25]{Cara memutuskannya tanpa berdebat}
\begin{isteps}
\item \{'h': 'Jalankan kumpulan uji Anda dengan kedua model', 'p': 'Model cepat dan model penalar, konfigurasi lain sama persis.'\}
\item \{'h': 'Catat tiga angka, bukan satu', 'p': 'Tugas yang selesai benar, biaya per tugas selesai, dan waktu dinding per tugas. Ketepatan saja akan selalu memenangkan model penalar.'\}
\item \{'h': 'Kalau selisih ketepatannya kecil, sudah selesai', 'p': 'Ambil yang murah. Selisih dua poin jarang sepadan dengan biaya sembilan kali lipat.'\}
\item \{'h': 'Kalau besar, jangan pukul rata', 'p': 'Cari langkah \textbf{mana} yang membaik. Biasanya satu atau dua, dan hanya itu yang perlu model penalar.'\}
\end{isteps}
\begin{iband}[signal]
Yang membuat langkah keempat mungkin: jejak per giliran. Tanpa itu, satu-satunya keputusan yang bisa diambil adalah keputusan pukul rata untuk seluruh sistem.
\end{iband}
\end{frame}

\framekicker{Pertanyaan yang selalu muncul}
\begin{frame}[shrink=25]{Kalau modelnya sudah bisa menalar, masih perlu gelung agen?}
{\fontsize{9.4}{12.6}\selectfont\color{fg2}Pertanyaan yang wajar, dan jawabannya bukan soal kecerdasan. Penalaran di dalam model dan gelung agen memecahkan \textbf{masalah yang berbeda}, dan hanya salah satunya bisa menyentuh dunia.\par}\vskip 4pt
\begin{minipage}[t]{0.485\textwidth}
{\fontsize{9.4}{12.6}\selectfont\color{fg2}\textbf{Yang bisa dilakukan penalaran sendiri}\par}\vskip 4pt
\begin{ibullets}
\item Memecah soal jadi langkah, di dalam satu panggilan
\item Memeriksa ulang hitungannya sendiri
\item Menyusun rencana sebelum menjawab
\end{ibullets}

\end{minipage}%
\hfill
\begin{minipage}[t]{0.485\textwidth}
{\fontsize{9.4}{12.6}\selectfont\color{fg2}\textbf{Yang tetap butuh gelung}\par}\vskip 4pt
\begin{ibullets}
\item Mengambil data yang belum ada di konteks
\item Melakukan sesuatu yang punya efek
\item Berhenti karena anggaran habis
\item Meninggalkan jejak yang bisa diaudit
\end{ibullets}

\end{minipage}%
\begin{iband}[signal]
Penalaran membuat \textbf{tiap langkah} lebih baik. Gelung yang membuat langkahnya ada sama sekali --- dan yang membuatnya bisa dihentikan, dibatasi, dan diperiksa. \hilite{Tidak ada jumlah penalaran yang menggantikan sebuah alat}
\end{iband}
\end{frame}

\coursesection{06}{Memakainya di dalam sistem}{Enam hal yang berubah begitu model penalar masuk ke gelung.}

\framekicker{Praktik}
\begin{frame}[shrink=25]{Penalaran dan singgahan prompt tidak saling membantu}
{\fontsize{9.4}{12.6}\selectfont\color{fg2}Singgahan prompt menghemat pada bagian \textbf{masukan} yang tidak berubah. Token penalaran adalah \textbf{keluaran}, dihasilkan baru tiap permintaan, dan tidak ada bagian darinya yang bisa disinggahkan.\par}\vskip 4pt
{\fontsize{9.4}{12.6}\selectfont\color{fg2}Akibatnya, dua pengungkit biaya terbesar Anda bekerja di sisi yang berbeda: singgahan menekan biaya masukan, anggaran berpikir menekan biaya keluaran. Keduanya perlu dipasang, dan salah satunya saja meninggalkan setengah penghematan di meja.\par}\vskip 4pt
\begin{iband}[signal]
Cara cepat melihat mana yang jadi masalah Anda: bandingkan token masukan dan keluaran per tugas. Kalau keluarannya besar, itu penalaran; kalau masukannya besar, itu riwayat yang tidak dipangkas.
\end{iband}
\end{frame}

\framekicker{Praktik}
\begin{frame}[shrink=25]{Penalaran dan keluaran berformat ketat saling menarik}
{\fontsize{9.4}{12.6}\selectfont\color{fg2}Model penalar diberi imbalan untuk sampai pada jawaban benar, bukan untuk patuh pada format. Pada tugas yang menuntut JSON dengan skema ketat, penalaran panjang kadang \textbf{menurunkan} kepatuhan format.\par}\vskip 4pt
\begin{isteps}
\item \{'h': 'Pisahkan berpikir dari mengeluarkan', 'p': 'Biarkan ia menalar, lalu minta panggilan alat sebagai langkah tersendiri. Dua panggilan yang masing-masing sederhana lebih andal daripada satu yang mengerjakan keduanya.'\}
\item \{'h': 'Validasi tetap wajib', 'p': 'Bab 2 sudah menyebutkan ini, dan penalaran tidak mengubahnya: skema menjamin bentuk, bukan kebenaran.'\}
\item \{'h': 'Catat laju penolakan skema per model', 'p': 'Kalau naik sesudah berpindah ke model penalar, Anda baru saja menemukan biaya tersembunyinya.'\}
\end{isteps}
\end{frame}

\framekicker{Praktik}
\begin{frame}[shrink=25]{Jejak penalaran bukan bukti alasan}
{\fontsize{9.4}{12.6}\selectfont\color{fg2}Menggoda sekali memperlakukan jejak penalaran sebagai penjelasan yang bisa diaudit: ia berbentuk alasan, berurutan, dan terbaca masuk akal.\par}\vskip 4pt
{\fontsize{9.4}{12.6}\selectfont\color{fg2}Tetapi tidak ada yang menjamin bahwa jejak itu \textbf{sebab} dari jawabannya. Model bisa sampai pada jawaban karena satu hal dan menuliskan urutan langkah yang berbeda --- dan itu tetap terbaca meyakinkan.\par}\vskip 4pt
\begin{iband}[signal]
Untuk keputusan yang diawasi, yang bisa diaudit adalah \textbf{panggilan alat dan hasilnya} --- itu peristiwa yang benar-benar terjadi. Jejak penalaran berguna untuk memperbaiki sistem; ia \hilite{bukan alat bukti}
\end{iband}
\end{frame}

\framekicker{Praktik}
\begin{frame}[shrink=25]{Menilainya: tiga angka, dan satu yang menyesatkan}
\begin{itable}{>{\raggedright\arraybackslash}p{0.2444\textwidth}>{\raggedright\arraybackslash}p{0.376\textwidth}>{\raggedright\arraybackslash}p{0.3196\textwidth}}
\thead{Ukuran} & \thead{Yang diberitahukannya} & \thead{Jebakannya} \\ \midrule
Tugas selesai benar & Yang sebenarnya Anda beli & Harus per tugas, bukan per panggilan \\
Biaya per tugas selesai & Harga sebenarnya & Harga per juta token menyembunyikan token tak terlihat \\
Waktu dinding per tugas & Apakah bisa dipakai orang & Rerata menyembunyikan ekor; lihat sebarannya \\
\textbf{Ketepatan per panggilan} & Terlihat paling bagus & \textbf{Selalu memenangkan model penalar}, dan tidak menjawab pertanyaan Anda \\
\end{itable}
{\fontsize{9.4}{12.6}\selectfont\color{fg2}Baris terakhir itu sebabnya perbandingan model sering berakhir dengan keputusan yang salah: ukuran yang paling mudah diambil adalah ukuran yang paling tidak relevan.\par}\vskip 4pt
\end{frame}

\framekicker{Praktik}
\begin{frame}[shrink=25]{Bentuk yang biasanya menang di produksi}
\begin{minipage}[t]{0.3253\textwidth}
\begin{icardgood}
\cardh{Model cepat di gelung}
Memilih alat, membaca hasil, memutuskan langkah berikutnya. Butuh kepatuhan skema, bukan penalaran panjang.
\end{icardgood}
\end{minipage}%
\hfill
\begin{minipage}[t]{0.3253\textwidth}
\begin{icardaccent}
\cardh{Model penalar di satu titik}
Langkah yang benar-benar sulit --- menyusun rencana, atau memutuskan kasus yang ambigu. Satu panggilan, dengan anggaran.
\end{icardaccent}
\end{minipage}%
\hfill
\begin{minipage}[t]{0.3253\textwidth}
\begin{icardgood}
\cardh{Kode untuk yang pasti}
Aritmetika, aturan kebijakan, validasi. Tidak ada model yang perlu menalar tentang hal yang bisa dihitung.
\end{icardgood}
\end{minipage}%
\begin{iband}[signal]
Perhatikan bahwa ini bentuk yang sama dengan demo kredit UMKM: \textbf{model klasik untuk angkanya, kode untuk kebijakannya, model bahasa untuk memutuskan apa yang dicari dan bagaimana menjelaskannya.}
\end{iband}
\end{frame}

\framekicker{Praktik}
\begin{frame}[shrink=25]{Kesalahan yang paling sering terjadi saat mengadopsinya}
\begin{minipage}[t]{0.494\textwidth}
\begin{icardbad}
\cardh{Memakainya untuk seluruh gelung}
Membayar penalaran pada enam langkah padahal hanya satu yang membutuhkannya. Ini kesalahan yang paling mahal dan paling mudah dihindari.
\end{icardbad}
\end{minipage}%
\hfill
\begin{minipage}[t]{0.494\textwidth}
\begin{icardbad}
\cardh{Tidak mengukur token tak terlihat}
Tagihan naik sembilan kali lipat dan tidak ada di dasbor mana pun, sebab yang dicatat hanya jawaban yang terlihat.
\end{icardbad}
\end{minipage}%


\vskip 4pt

\begin{minipage}[t]{0.494\textwidth}
\begin{icardbad}
\cardh{Menjadikan jejak penalaran sebagai alat bukti}
Terbaca meyakinkan, tidak dijamin sebagai sebab. Yang bisa diaudit adalah panggilan alat.
\end{icardbad}
\end{minipage}%
\hfill
\begin{minipage}[t]{0.494\textwidth}
\begin{icardbad}
\cardh{Membandingkan dengan ukuran yang salah}
Ketepatan per panggilan selalu memenangkan model penalar dan tidak pernah menjawab apakah sistemnya jadi lebih baik.
\end{icardbad}
\end{minipage}%
\end{frame}

\framekicker{Praktik}
\begin{frame}[shrink=25]{Titik awal yang masuk akal, sebelum diukur}
\begin{isteps}
\item \{'h': 'Mulai tanpa model penalar sama sekali', 'p': 'Bangun gelungnya, alatnya, dan kumpulan ujinya. Banyak sistem berhenti di sini dan sudah cukup.'\}
\item \{'h': 'Cari langkah yang paling sering salah', 'p': 'Dari jejaknya, bukan dari perasaan. Biasanya satu langkah menyumbang sebagian besar kegagalan.'\}
\item \{'h': 'Ganti langkah itu saja', 'p': 'Model penalar dengan anggaran berpikir yang dibatasi, hanya di titik itu.'\}
\item \{'h': 'Ukur tiga angka sebelum dan sesudah', 'p': 'Kalau biaya per tugas selesai naik lebih cepat daripada ketepatannya, kembalikan.'\}
\end{isteps}
\begin{iband}[signal]
Urutan ini terdengar lambat dan sebenarnya paling cepat: ia menghindari satu-satunya kesalahan yang mahal, yaitu memakai model penalar di mana-mana lalu mencari tahu belakangan bahwa hanya satu langkah yang membutuhkannya.
\end{iband}
\end{frame}

\framekicker{Praktik}
\begin{frame}[shrink=25]{Satu hal yang tidak diperbaiki oleh penalaran, sekuat apa pun}
{\fontsize{9.4}{12.6}\selectfont\color{fg2}Model penalar tetap tidak tahu apa yang tidak diketahuinya. Penalaran yang panjang di atas fakta yang salah menghasilkan kesimpulan yang salah \textbf{dengan lebih banyak dukungan} --- bukan lebih sedikit.\par}\vskip 4pt
{\fontsize{9.4}{12.6}\selectfont\color{fg2}Itu sebabnya bab berikutnya bukan tentang membuat model berpikir lebih baik, melainkan tentang \textbf{memberinya bahan yang benar}: apa yang diingat, apa yang dicari lagi, dan apa yang dibuang.\par}\vskip 4pt
\begin{iband}[signal]
Urutan yang benar hampir selalu: perbaiki bahannya dulu, baru pertimbangkan menaikkan kemampuan menalarnya. Terbalik dari yang biasanya dicoba orang.
\end{iband}
\end{frame}

\framekicker{Penutup}
\begin{frame}[shrink=25]{Yang dibawa pulang dari bab ini}
\begin{isteps}
\item \{'h': 'Kualitas kini bisa dibeli saat pemakaian', 'p': 'Dan karena itu ia jadi keputusan per permintaan, bukan sifat tetap sistem.'\}
\item \{'h': 'Yang dibeli melandai, yang dibayar tidak', 'p': '1 $\rightarrow$ 5 contoh: +8,3 poin dengan biaya 5$\times$. 5 $\rightarrow$ 21: +14,3 poin dengan biaya 21$\times$. Dan itu pun batas atas.'\}
\item \{'h': 'Pemeriksa mengalahkan contoh tambahan', 'p': 'Asal memeriksa lebih mudah daripada menjawab --- dan pemeriksa terbaik biasanya kode biasa.'\}
\item \{'h': 'Sembilan dari sepuluh token penalaran tidak terlihat', 'p': 'Dan di dalam gelung agen, ia dikalikan jumlah langkah.'\}
\item \{'h': 'Penalaran panjang punya arah gagalnya sendiri', 'p': 'Jawaban salah yang datang dengan alasan panjang lebih sulit ditolak daripada tebakan singkat.'\}
\end{isteps}
\end{frame}
\note{Pertanyaan penutup yang bagus untuk kelas: di sistem kalian, langkah mana yang benar-benar butuh penalaran? Hampir selalu jawabannya satu atau dua langkah, bukan seluruh gelung.}

\end{document}
